\documentclass[12pt]{article}
\usepackage[utf8]{inputenc}
\usepackage[T1]{fontenc}
\usepackage{hyperref}
\title{Kunbarlang Language Resources for K-12 Curriculum Development}
\author{Compiled by Claude Sonnet 5.5 (claude-sonnet-5-5) for Andr\'as Kornai}
\date{}
\begin{document}
\maketitle

\section{Introduction}
Kunbarlang (Kun-barlang, Gunbarlang; ISO 639-3 wlg) is a Gunwinyguan language of central coastal Arnhem Land, Northern Territory, Australia. Sources conflict on its status. ELDP records describe about 60 remaining speakers in one collection and fewer than 50 in a later one, while Wikipedia lists it as extinct by 2016. It should be treated as severely endangered. The aim here is heritage and revitalization materials.

\section{Primary Educational Resources}
\begin{itemize}
\item A children's picture dictionary of Kun-barlang, with main entries in Kunbarlang and subordinate entries for the Kun-kurdu-wala variety. It was published by the Warruwi Literature Production Centre and compiled by Ngal-kangila Barbara Warramamkudj, Ngal-kangila Betty Galawarridja and Na-kangila Paul Naragoidj. It is held as a PDF in the Territory Stories repository (licence CC BY-NC-ND 4.0) and catalogued by AustLit. I do not have a confirmed link for that specific item, so no URL is given here.
\item A Kunbarlang children's text first published in 1984 by the Maningrida Literature Production Centre, also held by Territory Stories. Its spelling system is almost identical to that used for Kuninjku. No confirmed link for that specific item is given.
\end{itemize}

\section{Linguistic and Reference Materials}
No full published grammar or dictionary was found in search results.

\section{Religious and Community Resources}
No scripture or community organisation resources were found in search results.

\section{Digital and Audio Resources}
\begin{itemize}
\item ELDP collection 0384, ``Varieties of Kun-barlang, a language on the edge'' (depositors Isabel O'Keeffe and Ruth Singer), a community- and research-accessible corpus: \url{https://eldp.access.preservica.com/dk0384/}
\item ELDP collection 0400, ``Comprehensive pan-varietal, ethnobiological, anthropological record of Kun-barlang'' (depositor Isabel O'Keeffe): \url{https://eldp.access.preservica.com/dk0400/}
\end{itemize}

\section{Curriculum Development Context}
The documentation corpora and the children's picture dictionary are the best starting points for heritage materials. Any work should be done with the remaining speakers and their communities, who hold the rights to these materials.
\end{document}